%%
%% This is file `lexample.tex',
%% Sample file for siam macros for use with LaTeX 2e
%%
%% By Paul Duggan for the Society for Industrial and Applied
%% Mathematics.
%%
%% October 1, 1995
%%
%% Version 1.0
%%
%% You are not allowed to change this file.
%%
%% You are allowed to distribute this file under the condition that
%% it is distributed together with all of the files in the siam macro
%% distribution. These are:
%%
%%  siamltex.cls (main LaTeX macro file for SIAM)
%%  siamltex.sty (includes siamltex.cls for compatibility mode)
%%  siam10.clo   (size option for 10pt papers)
%%  subeqn.clo   (allows equation numbners with lettered subelements)
%%  siam.bst     (bibliographic style file for BibTeX)
%%  docultex.tex (documentation file)
%%  lexample.tex (this file)
%%
%% If you receive only some of these files from someone, complain!
%%
%% You are NOT ALLOWED to distribute this file alone. You are NOT
%% ALLOWED to take money for the distribution or use of either this
%% file or a changed version, except for a nominal charge for copying
%% etc.
%% \CharacterTable
%%  {Upper-case    \A\B\C\D\E\F\G\H\I\J\K\L\M\N\O\P\Q\R\S\T\U\V\W\X\Y\Z
%%   Lower-case    \a\b\c\d\e\f\g\h\i\j\k\l\m\n\o\p\q\r\s\t\u\v\w\x\y\z
%%   Digits        \0\1\2\3\4\5\6\7\8\9
%%   Exclamation   \!     Double quote  \"     Hash (number) \#
%%   Dollar        \$     Percent       \%     Ampersand     \&
%%   Acute accent  \'     Left paren    \(     Right paren   \)
%%   Asterisk      \*     Plus          \+     Comma         \,
%%   Minus         \-     Point         \.     Solidus       \/
%%   Colon         \:     Semicolon     \;     Less than     \<
%%   Equals        \=     Greater than  \>     Question mark \?
%%   Commercial at \@     Left bracket  \[     Backslash     \\
%%   Right bracket \]     Circumflex    \^     Underscore    \_
%%   Grave accent  \`     Left brace    \{     Vertical bar  \|
%%   Right brace   \}     Tilde         \~}


\documentclass[final]{siamltex}

% siamltex.cls does not load the AMS packages; without them \text, cases,
% align*, \tag, \eqref and \mathbb fail and the output is garbled.
\usepackage{amsmath,amssymb}
% \bysame and \MR come from amsart; define them for the bibliography.
\providecommand{\bysame}{\leavevmode\hbox to3em{\hrulefill}\thinspace}
\providecommand{\MR}[1]{MR #1}
% Unnumbered remark (not provided by siamltex.cls).
\makeatletter
\@ifundefined{rem}{\newenvironment{rem}{\par\noindent{\em Remark.}\ }{\par}}{}
\makeatother

\newcommand{\g}{\gamma}
\newcommand{\ep}{\varepsilon}
\newcommand{\ve}{\varepsilon}
\newcommand{\epj}{\varepsilon_j}
\newcommand{\epn}{\varepsilon_n}
\newcommand{\epnk}{\varepsilon_{n_k}}
\newcommand{\psin}{\psi^n}
\newcommand{\uep}{u^\varepsilon}
\newcommand{\uepj}{u^{\varepsilon_j}}
\newcommand{\uepk}{u^{\varepsilon_k}}
\newcommand{\vepj}{v^{\varepsilon_j}}
\newcommand{\uepnk}{u^{\varepsilon_{n_k}}}
\newcommand{\wep}{w^\varepsilon}
\newcommand{\zep}{z^\varepsilon}
\newcommand{\wepj}{w^{\varepsilon_j}}
\newcommand{\zepj}{z^{\varepsilon_j}}
\newcommand{\pep}{P_\varepsilon}
\newcommand{\cald}{\mathcal{D}}
\newcommand{\aej}{{\mathcal{A}}_{\varepsilon_j}}
\newcommand{\ub}{{\bar u}^{\varepsilon}}
\newcommand{\vb}{{\bar v}^{\varepsilon}}
\newcommand{\wb}{{\bar w}^{\varepsilon}}
\newcommand{\iy}{\infty}
\newcommand{\yep}{Y^\varepsilon}
\newcommand{\vep}{v^\varepsilon}
\newcommand{\uepn}{u^{\varepsilon_n}}
\newcommand{\xe}{x_{\varepsilon}}
\newcommand{\te}{t_{\varepsilon}}
\newcommand{\ubj}{{\bar u}^{\varepsilon_j}}
\newcommand{\ye}{y_{\varepsilon}}
\newcommand{\wu}{\widehat u}
\newcommand{\pa}{\partial}
\newcommand{\R}{\mathbb{R}}
\newcommand{\Rr}[1]{\mathbb{R}^{#1}}
\newcommand{\ubar}{\overline u}
\newcommand{\wt}{\widetilde}
\newcommand{\al}{\alpha}
\newcommand{\m}{\sqrt{2M(x_0,t_0)}}
\newcommand{\mxt}{\sqrt{2M(x,t)}}

\title{Uniqueness of limit
solutions to a free boundary problem from
combustion\thanks{Published in SIAM J. Math. Anal. {\bf 36} (2004), no.~1, 172--185.
Partially supported by grant BID1201/OC-AR PICT
03-00000-05009 and CONICET grant PIP0660/98. The second author is
a member of CONICET (Consejo Nacional de Investigaciones
Cient\'{\i }ficas y T\'{e}cnicas of Argentina).}}

% The thanks line in the title should be filled in if there is
% any support acknowledgement for the overall work to be included
% This \thanks is also used for the received by date info, but
% authors are not expected to provide this.

\author{Juli\'{a}n Fern\'{a}ndez Bonder\thanks{Departamento de Matem\'{a}tica, Facultad de Ciencias Exactas
y Naturales, Universidad de Buenos Aires, (1428), Buenos Aires,
Argentina ({\tt jfbonder@dm.uba.ar}).}
        \and Noemi Wolanski\thanks{Departamento de Matem\'{a}tica, Facultad de Ciencias Exactas
y Naturales, Universidad de Buenos Aires, (1428), Buenos Aires,
Argentina ({\tt wolanski@dm.uba.ar}).}}

\begin{document}

\maketitle

\begin{abstract}
We investigate the uniqueness of limit solutions for a free
boundary problem in heat propagation that appears as a limit of a
parabolic system that arises in flame propagation.
\end{abstract}

\begin{keywords}
Free-boundary problem, combustion, heat equation, uniqueness,
classical solution, limit solution
\end{keywords}

\begin{AMS}
35K05, 35K60, 80A25
\end{AMS}

\pagestyle{myheadings}
\thispagestyle{plain}
\markboth{J. Fern\'{a}ndez Bonder and N. Wolanski}{Uniqueness of limit solutions to a FBP from combustion}


\section{Introduction}
\setcounter{equation}{0}

In this paper we consider the following problem arising in combustion theory

\begin{equation}
\label{1.1}
\left\{\begin{array}{rcll}
\Delta u^{\ve} - u^{\ve}_{t} & = & Y^{\ve} f_{\ve}(u^{\ve})
& \mbox{in } {\mathcal D},\\
\Delta Y^{\ve} - Y^{\ve}_{t} & = & Y^{\ve} f_{\ve}(u^{\ve})
& \mbox{in } {\mathcal D},\\
\end{array}\right.
\end{equation}
where ${\mathcal D}\subset\R^{N+1}$.

This model appears in combustion theory in the analysis of the
propagation of curved flames. It is derived in the framework of
the theory of equidiffusional premixed flames analyzed in the
relevant limit of high activation energy for Lewis number 1. In
this application, $Y^{\ve}$ represents the fraction of some
reactant (and hence it is assumed to be nonnegative), and
$u^{\ve}$ is minus the temperature (more precisely,
$u^{\ve}=\lambda(T_f-T^{\ve})$ where $T_f$ is the flame
temperature  and $\lambda$ is a normalization factor). Observe
that the term $Y^{\ve}f_{\ve}(u^{\ve})$ acts as an absorption term
in the equation (\ref{1.1}). Since $T^{\ve} = T_f -
(u^{\ve}/\lambda)$, it is in fact a reaction term for the
temperature. In the flame model, such a term represents the effect
of the exothermic chemical reaction and $f$ has accordingly a
number of properties: it is a nonnegative Lipschitz continuous
function which is positive in an interval $(-\infty, \ve)$ and
vanishes otherwise (i.e., reaction occurs only when
$T>T_f-\frac\ve\lambda$). The parameter $\ve$ is essentially the
inverse of the activation energy of the chemical reaction. For the
sake of simplicity we will assume that
$f_{\ve}(s)=\frac{1}{\ve^2}f(\frac{s}{\ve}),$ where $f$ is a
Lipschitz continuous function with $f(s)>0$ if $s<1$ and $f(s)=0$ if $s\ge1$.

For the derivation of the model, we cite \cite{BuL}.

Here we are interested in high activation energy limits (i.e. $\ve\to 0$).
These limits are currently the subject of active investigation, specially
in the case $u^{\ve}=Y^{\ve}$. This is a natural assumption in the case of travelling waves.
%%%%%%%%%%%%%%%%%%%%%%%%%%%%%%%%%%
%%%%%%%%%%%%%%%%%%%%%%%%%%

In a previous paper \cite{FBW} we have studied this problem in the case in which the initial values for $\uep$ and $\yep$
-- both  converging to the same function $u_0$ -- satisfy the condition
\begin{equation}
\label{condicion_inicial}
\frac{\yep_0(x)-\uep_0(x)}\ep\to w_0(x)\quad\mbox{uniformly  in } \R^N
\end{equation}
with $w_0>-1$.

Problem (\ref{1.1}) reduces to a single equation, namely
\begin{equation*}
\Delta\uep-\uep_t=(\uep+\wep)f_\ep(\uep),
\tag{$\pep$}
\end{equation*}
where the function $\wep(x,t)$ is the solution of the heat equation with initial
datum $\yep_0(x)-\uep_0(x)$. Observe that $\uep+\wep=\yep\ge0$.

By
(\ref{condicion_inicial}) there exists the limit
\begin{equation}
\label{def_w_0}
\lim_{\ep\to0}\frac{\wep(x,t)}\ep=w_0(x,t)
\end{equation}
and $w_0(x,t)$ is the solution of the heat equation with initial datum $w_0(x)$.

In this way, at least
formally, the reaction term converges to a delta function and a free
boundary problem appears. In fact, we have proved in \cite{FBW}
that every sequence of uniformly bounded solutions to (\ref{1.1}), \{$\uepn\}$, with $\epn\to0$ has  a subsequence
\{$\uepnk$\} converging to a  limit function $u\ge0$ which is a solution of the
following free boundary problem
\begin{equation*}
\left\{
\begin{aligned}
&\Delta u - u_t  =  0  \qquad\ \quad\ \ \ \mbox{in }  \{u>0\}\\
&|\nabla u^+|  =  \sqrt{2M(x,t)} \quad \mbox{on } \partial\{u>0\}
\end{aligned}
\right.\tag{$P$}
\end{equation*}
where $M(x,t)=\int_{-w_0(x,t)}^1\big(s+w_0(x,t)\big)f(s)ds$.



We see that the free boundary condition strongly depends on the approximation $\uep_0$, $\yep_0$ of the initial
datum $u_0$. In particular, the limit function $u$ is different for different approximations of the initial datum $u_0$.

It is therefore natural to wonder whether the only condition that determines the limit function $u$  is condition (\ref{condicion_inicial}).

The purpose of  this paper is to prove that this is indeed the case, at least under some monotonicity assumption on the initial value $u_0$. This monotonicity assumption is similar to that used  to prove uniqueness of the limit for the case $\uep=\yep$ in \cite{CP}.

In fact, we follow here some of the ideas of \cite{CP} which are based on the fact that any limit function is a supersolution to $(P)$. This is still true in our case. Unfortunately the simple construction in \cite{CP} of supersolutions of ($\pep$)
that approximate a strict classical  supersolution of ($P$), when $\wep=0$, does not work in the general case unless one
asks for a lot of complementary conditions on the reaction function $f$.

Therefore, we follow here the construction done in \cite{LVW1}. The proof that this
construction works is based on
blow up of the constructed functions. This technique was already seen to work very well for ($\pep$)  -- under the condition
(\ref{condicion_inicial}) -- in \cite{FBW}.

\medskip
Our result can be summarized as saying that, {\em under suitable assumptions on the domain and on the initial datum $u_0$,
there exists at most one limit solution to the free boundary problem $(P)$ with nonvanishing gradient near its free boundary, as long as the approximate initial data -- converging uniformly to $u_0$ with supports that converge to
the support of $u_0$ --  satisfy (\ref{condicion_inicial})}.

Moreover, under the same geometric assumptions, {\em if there exists a classical solution to $(P)$, this is the only
limit of solutions to ($\pep$) with initial data satisfying the conditions above. In particular, it is the
only classical solution to $(P)$}.

\medskip

As already stated, in the case $\uep=\yep$, uniqueness results for limit solutions under
geometric hypotheses similar to the ones made here can be found in \cite{CP}. Also in
\cite{LVW1} the authors study the uniqueness and agreement between different concepts of
solutions of problem ($P$) (again in the case $\uep=\yep$) under the assumption of the existence
of a classical solution and under different geometric assumptions. See also \cite{LVW2} for a similar result in the two-phase case.






\subsection*{Notation}
Throughout the paper $N$ will denote  the spatial dimension.
In addition, the following notation will be used:

For any $x_{0}\in \ {\mathbb{R}}^{N}$, $t_{0}\in \R $ and $\tau >0$,
$B_{\tau }(x_{0}):=\{x\in {\mathbb{R}}^{N} /|x-x_{0}|<\tau \}$ and
$B_{\tau }(x_{0},t_{0}):=\{(x,t)\in {\mathbb{R}}^{N+1} /\emph{}|x-x_{0}|^{2}+
|t-t_{0}|^{2}<{\tau }^{2}\}$.

When necessary, we will denote points in $\mathbb{R}^{N}$ by $x=(x_{1},x')$,
with
$x'\in \mathbb{R}^{N-1}$.
Given a function $v$, we will denote $v^{+}=\text{\rm max}(v,0)$.

The symbols $\Delta $ and $\nabla $ will denote the corresponding
operators in the space variables; the symbol ${\partial }_{p}$ applied
to a domain will denote parabolic boundary.



Finally, we will say that $u$ is supercaloric if $\Delta u - u_t\le 0$, and $u$
is subcaloric if $\Delta u - u_t \ge 0$.






\subsection*{Outline of the paper}
An outline of the contents is as follows.  In Section 2 we  give precise
definitions of classical sub- and supersolutions and prove a
comparison result for problem ($P$) (Lemma \ref{compar}).
In Section 3 we state some auxiliary results.   In
Section  4  we prove that a  strict classical  supersolution  to  problem ($P$) is
the uniform  limit of a  family of supersolutions  to  problem ($\pep$) (Theorem \ref{approx}) and as a consequence
we obtain the boundedness of the support for limit solutions in the geometry under consideration
(Proposition \ref{acotada}).
Finally, in Section 5 we prove our main result (Theorem \ref{unicidad}).
We discuss in a final section (Section 6) the results proved in the paper as well as other possible geometries that can be considered.

\section{Preliminaries}
\setcounter{equation}{0}

Following \cite{CP} we will define what we will understand by a classical
supersolution of problem ($P$). Note that the meaning of {\em classical} here differs from
the usual one since we are not assuming that the function is $C^1$ up to the free boundary
or that the free boundary is $C^1$.

\begin{definition}
\label{defsuper}
A continuous nonnegative function $u$ in
$\overline{Q}_T = \R^N \times [0,T],\ T>0$, is called
a classical supersolution of ($P$) if $u\in C^1(\{u>0\})$ and
\begin{itemize}
\item[{(i)}] \ $\Delta u - u_t \le 0 \qquad \qquad \text{in } \Omega=\{ u>0 \}$;
\item[{(ii)}] \ $\limsup_{\Omega\ni (y,s)\to (x,t)}|\nabla u(y,s)|\le \sqrt{2M(x,t)}$ for
every $(x,t)\in \partial\Omega\cap Q_T$;
%\item[{(iii)}] \ $\partial u/\partial\eta = 0 \qquad \qquad \text{in } \partial_N Q_T$
\item[{(iii)}] \  $u(\cdot,0)\ge u_0$.
\end{itemize}
Respectively, $u$ is a classical subsolution of ($P$) if conditions (i), (ii) and
(iii) are satisfied with reversed inequalities and lim inf instead of lim sup in (ii).

A function $u$ is a classical solution of ($P$)  if it is both
a classical subsolution and a classical supersolution of ($P$).
%, $u\in C^1(\overline{\{u>0\}})$ and the free boundary ${\partial\{u>0\}\cap Q_T}$
%is a $C^1$ surface.

Next, a classical supersolution $u$ of ($P$) is a strict classical supersolution of ($P$) if there is a
$\delta>0$ such that the stronger inequalities
\begin{itemize}
\item[{(ii')}]  \ $\lim\sup_{\Omega\ni (y,s)\to (x,t)}|\nabla u(y,s)|\le \sqrt{2M(x,t)-\delta}$ for
every $(x,t)\in \partial\Omega\cap Q_T$;
\item[{(iii')}] \  $u(\cdot,0)\ge u_0+\delta$ on $\Omega_0=\{ u_0>0 \}$
\end{itemize}
hold. Analogously a strict classical subsolution is defined.
\end{definition}


As a consequence of the results in \cite{FBW}, one can check that every limit
solution $u=\lim_{j\to \infty} \uepj$ of ($P$) is a
classical supersolution in the sense of Definition \ref{defsuper}. In fact,
\begin{proposition}
\label{rem1}
Let $\uepj$ be solutions to ($P_{\epj}$) -- with $w^{\ep_j}$ satisfying (\ref{def_w_0}) and $w_0>-1$ -- such that $\uepj \to u$ uniformly on
compact sets and $\epj\to 0$. Assume that the initial datum $u_0$ is Lipschitz
continuous and that the approximations of the initial datum verify
$|u_{0}^\ep(x)|, |\nabla u_0^\ep(x)|\le C$ and $u_0^\ep\in C^1(\overline{\{u_0^\ep>0\}})$.
Then $u$ is a classical supersolution of ($P$).
\end{proposition}

\begin{proof}
We have to verify conditions (i)-(iii) of Definition \ref{defsuper}.

{}From our assumptions on the initial datum $u_0$, by Proposition 5.2.1 of \cite{L}, we have that
$\uep\to u$ uniformly on compact sets of $\overline{Q_T}$ so that
$u$ is continuous up to $t=0$ and (iii) holds.

Now, (i) is proved in \cite{FBW}.

Finally, (ii) is a straightforward modification of Theorem 6.1 of \cite{CLW1} using Lemmas 2.1, 2.2 and 2.3
of \cite{FBW} instead of Lemma 3.2 and Propositions 5.2 and 5.3 of \cite{CLW1} respectively.
\end{proof}

\medskip



\medskip

Let us suppose that the initial datum $u_0$ of problem ($P$) is
starshaped with respect to a point $x_0$, that we always assume to be 0, in the following sense:
For  every $\lambda\in (0,1)$ and $x\in\R^N$,
\begin{equation}
\label{S}
u_0(\lambda x)\ge u_0(x),  \qquad
 \lambda \Omega_0 \subset \subset \Omega_0,
\end{equation}
where $\Omega_0 = \{ u_0 > 0 \}$.

Also, assume that
\begin{equation}
\label{mono_w}
w_0(\lambda x,0)\le w_0(x,0)\quad\mbox{ if }x\in\R^N\ ,\  0<\lambda<1.
\end{equation}

\medskip

Let $u$ be a  classical supersolution of ($P$). Let $\lambda$ and $\lambda'$ be two real numbers
with $0<\lambda<\lambda'<1$. Define
\begin{equation}
\label{rescalada}
u_{\lambda}(x,t) = \frac{1}{\lambda'} u(\lambda x, \lambda^2 t)
\end{equation}
in $Q_{T/\lambda^2}$. The rescaling is taken so that $u_\lambda$ is a supersolution of the
heat equation in
\begin{equation}
\label{omel}
\Omega_{\lambda} = \{ (x,t):\ (\lambda x,\lambda^2 t)\in \Omega \}.
\end{equation}
Moreover, the fact that $0<\lambda<\lambda'<1$ makes $u_\lambda$ a strict classical
 supersolution of ($P$).

In fact, let us first see that
\begin{equation*}
M(\lambda x,\lambda^2t)\le M(x,t)\qquad\mbox{if }0<\lambda<1.
\end{equation*}
This is a consequence of the fact that the function
$$
a\longrightarrow \int_{-a}^1(s+a)f(s)\,ds
$$
is nondecreasing and
\begin{equation}
\label{mono_w2}
w_0(\lambda x,\lambda^2 t)\le w_0(x,t)\qquad\mbox{if }0<\lambda<1.
\end{equation}
In fact, the function $w_\lambda(x,t)=w_0(\lambda x,\lambda^2 t)$ is caloric and $w_\lambda(x,0)\le w_0(x,0)$ if
$0<\lambda<1$
by hypothesis. Thus, by the comparison principle, $w_\lambda(x,t)\le w_0(x,t)$ in $\R^N\times(0,T)$.

Now, let $(x_0,t_0)\in\partial\{u_\lambda>0\}$. Then,
\begin{equation*}
\begin{aligned}
\limsup_{\Omega_\lambda\ni (x,t)\to (x_0,t_0)}|\nabla u_\lambda(x,t)|&=
\limsup_{\Omega\ni (\lambda x,\lambda^2 t)\to (\lambda x_0,\lambda^2 t_0)}|\frac{\lambda}{\lambda'}\nabla u(\lambda x,\lambda^2 t)|\\
&\le\frac{\lambda}{\lambda'}\sqrt{2M(\lambda x_0,\lambda^2 t_0)}\le \sqrt{2M(x_0,t_0)}-\Big(1-\frac\lambda{\lambda'}\Big)\sqrt{2M_0},
\end{aligned}
\end{equation*}
where $0<M_0<M(x,t)$ in $\R^N\times(0,T)$.

\medskip

On the other hand, since $\lambda \Omega_0\subset\subset\Omega_0$, there holds that
$$
u_0(\lambda x)\ge\gamma>0\qquad\mbox{if }x\in \Omega_0.
$$
Thus, for $x\in \Omega_0$,
\begin{align*}
u_\lambda(x,0)&=\frac1{\lambda'}u_0(\lambda x)=u_0(\lambda x)+\Big(\frac1{\lambda'}-1\Big)u_0(\lambda x)\\
&\ge u_0(x)+\Big(\frac1{\lambda'}-1\Big)\gamma.
\end{align*}

\medskip

The following comparison lemma for problem ($P$) can be proved as Lemma 2.4 in \cite{CP}. We omit the proof.

\begin{lemma}
\label{compar}
Let $u_0$ satisfy \eqref{S} and $w_0$ satisfy (\ref{mono_w}). Then every classical subsolution of ($P$) with bounded support, is smaller than every classical
supersolution of ($P$). i.e. if $u'$ is a classical subsolution such that $\Omega'$ is bounded and $u$ is a classical supersolution then
$$\Omega'\subset \Omega \quad \mbox{and} \quad u'\le u,$$
where $\Omega' = \{ u'>0 \}$ and $\Omega = \{ u>0 \}$.
\end{lemma}


\section{Auxiliary results}
\setcounter{equation}{0}


This section contains results on the following problem:
\begin{equation*}\Delta u - u_{t}=(u + \omega_0)f(u),
\tag{$P_{0}$}
\end{equation*}
where the function $f $ is as in Section 1 and $\omega_0$ is a constant, $\omega_0>-1$.
The results will be used
in the next sections where $(P_{0})$ appears as a blow-up limit. The proofs are very similar to those of Lemmas 4.1, 4.3 and 4.4 in \cite{LVW1}.
We leave the details to the reader.

\begin{lemma}\label{4.1}
Let $a,b\geq 0$ and let $\psi $ be the classical solution to
\begin{equation}\label{eq4.1}\begin{split}
\psi _{ss}&=(\psi + \omega_0)f (\psi )\quad \text{for }s>0,\\
\psi (0)&=a, \quad \psi _{s}(0)=-\sqrt {2b}.
\end{split}\end{equation}
Let $B(\tau )=\int _{-\omega_0}^{\tau } (\rho+\omega_0)f (\rho )\,d\rho $.
\begin{align}
\label{4.2}& \mbox{If } b=0 \mbox{ and } a\in \{-\omega_0\}\cup [1,+\infty ), \mbox{ then }\psi \equiv a.\\
\label{4.3}& \mbox{If $b=0$ and $a\in (-\omega_0,1)$, then $\lim _{s\to +\infty } \psi (s)=+\infty $.}\\
\label{4.4}& \mbox{If $b\in \bigl (0,\,B(a)\bigr )$,
then $\lim _{s\to +\infty } \psi (s)=+\infty $.}\\
\label{4.5}& \mbox{If $0<b=B(a)$, then $\psi _{s}<0$ and $\lim _{s\to +\infty } \psi (s)=-\omega_0$.}\\
\label{4.6}& \mbox{If $b\in \bigl (B(a),\, +\infty \bigr )$, then $\psi _{s}<0$ and
$\lim _{s\to +\infty } \psi (s)=-\infty $.}
\end{align}
\end{lemma}


\begin{lemma}\label{lema4.3}
Let $B(\tau)$ be as in the previous Lemma and let
$\mathcal{R}_{\gamma }=\bigl \{(x,t)\in {\mathbb{R}}^{N+1}\,/x_{1}>0\,,\,-\infty <t\le \gamma \bigr \}$,
 $0\leq \theta <1+\omega_0$ and let
$U\in C^{2+\alpha,1+\frac{\alpha}{2}}(\overline{\mathcal{R}_{\g }})$
be such that
\begin{alignat*}{2}
\Delta U-U_{t}&=(U + \omega_0)f (U) &\quad &\text{in }\mathcal{R}_{\g },\\
U&=1-\theta & \quad &\text{on } \{x_{1}=0\},\\
-\omega_0\leq &U\leq 1-\theta& \quad &\text{in }\overline{\mathcal{R}_{\g }}.
\end{alignat*}

1) If $\theta =0$, then $|\nabla U|\leq \sqrt {2B(1)}$ on $\{x_{1}=0\}$.

2) If $0<\theta <1+\omega_0$ and $0<\sigma <B(1)$ are such that
$\int _{-\omega_0}^{1-\theta } (\rho + \omega_0)f (\rho )\,d\rho =B(1)-\sigma $,
then $|\nabla U|= \sqrt {2(B(1)-\sigma )}$ on $\{x_{1}=0\}$.
\end{lemma}

Finally, we state a compactness result.

\begin{lemma} \label{lema4.4}
Let $\epj $, $\g _{\epj }$ and $\tau _{\epj }$ be sequences such  that
$\epj >0$, $\epj \to 0$, $\g _{\epj }>0$, $\g _{\epj }\to \g $, with $0\le \g \le +\infty $,
$\tau _{\epj }>0$, $\tau _{\epj }\to \tau $ with
$0\leq \tau \leq +\infty $, and such that $\tau <+\infty $ implies that $\g =+
\infty $. Assume that $w^{\ep_j}/\ep_j$ converge to $w_0$ uniformly in compact sets of $\R^N\times[0,T]$.
Let $\rho >0$ and
\begin{equation*}\aej =\left \{ (x,t)\,\,/\,\, |x|<\frac{\rho }{\epj }\, ,
\,\,-\min (\tau _{\epj },\frac{\rho ^{2}}{\epj ^{2}})<t<\min (
\g _{\epj },\frac{\rho ^{2}}{\epj ^{2}})\right \}.\end{equation*}
Let $(x_0,t_0)\in \R^N\times[0,T]$.
Assume that $0\leq \theta <1 + w_0(x_0,t_0)$ and let $\ubj $ be weak solutions to
\begin{equation*}
\begin{split}
& \Delta \ubj - \ubj_t  = \left(\ubj + \frac{\wepj(\epj x + x_{\epj}, \epj^2 t + t_{\epj})}{\epj}\right)f(\ubj )
\text{ in }\{x_{1}>\bar h_{\epj }(x',t)\}\cap \aej ,\\
& \ubj =1-\theta \quad \text{on }\{x_{1}=\bar h_{\epj }(x',t)\}\cap \aej ,\\
-&\frac{\wepj(\epj x + x_{\epj}, \epj^2 t + t_{\epj})}{\epj}\leq \ubj \leq 1-\theta \quad \text{in }\{x_{1}\geq \bar h_{\epj
}(x',t)\}\cap \overline{\aej },
\end{split}
\end{equation*}
where $(x_{\epj},t_{\epj})\to (x_0,t_0)$,
with $\ubj \in C(\{x_{1}\ge \bar h_{\epj }(x',t)\}\cap \overline{\aej })$, and
$\nabla \ubj \in L^{2}$.
Here $\bar h_{\epj }$ are continuous functions such that $\bar h_{\epj }(0,0)=0$
with $\bar h_{\epj }\to 0$ uniformly on compact subsets of
${\mathbb{R}}^{N-1}\times (-\tau ,\g )$.  Moreover, we assume that
$\| \bar h_{\epj }\|_{C^{1}(K)} +
\| \nabla _{x'} \bar h_{\epj }\|_{C^{\alpha,\frac{\alpha}{2}}(K)}$ are uniformly
bounded, for every compact set $K\subset \mathbb{R}^{N-1}\times (-\tau ,\g )$.

Then, there exists a function $\bar u$ such that, for a subsequence,
\begin{gather*}
\bar u\in C^{2+\alpha ,1+ \frac{\alpha }{2}}\bigl (\{x_{1}\geq 0,\,\g >t>-\tau
\}\bigr ),\\
\ubj \to \bar u\quad \text{uniformly on compact subsets of }
\{x_{1}>0,\,\g >t>-\tau \},\\
\begin{array}{rll}
&\Delta \bar u-\bar u_{t}=(\bar u + w_0(x_0,t_0))f(\bar u)&\quad \text{in
}\{x_{1}>0,\,\g >t>-\tau \},\\
&\bar u=1-\theta &\quad \text{on }\{x_{1}=0,\,\g >t>-\tau \},\\
-\hskip-10pt&w_0(x_0,t_0)\leq \bar u\leq 1-\theta &\quad \text{in }\{x_{1}\geq 0,\,\g >t>-\tau \}.
\end{array}
\end{gather*}

If $\gamma <+\infty $, we require, in addition, that
\[ \|\bar h_{\epj }(x',t+\g _{\epj }-\g )\|_{C^{1}(K)} +
\|\nabla _{x'}\bar h_{\epj }(x',t+\g _{\epj }-\g)\|_{C^{\alpha,\frac{\alpha}{2}}(K)}\]
be uniformly bounded for every compact set $K\subset \mathbb{R}^{N-1}\times (-\infty ,\g ]$. And we deduce that
\begin{equation*}\ubar \in C^{2+\alpha ,1+ \frac{\alpha }{2}}\bigl (\{x_{1}\geq 0,\,t\le \g \}\bigr ).
\end{equation*}

If $\tau <+\infty $,  we let
\begin{equation*}\mathcal{B}_{\epj }=\left \{ x\,\,/\,\, |x|<\frac{\rho }{\epj }\, ,
\,\,x_{1}>\bar h_{\epj }(x',-\tau _{\epj })\right \},\end{equation*}
and we require, in addition, that  for every $R>0$,
\begin{equation*}\|\ubj (x,-\tau _{\epj })\|_{C^{\alpha}\left (
\overline{ \mathcal{B}}_{\epj }\cap \overline{B}_{R} (0)
\right )}\leq C_{R},\end{equation*}
and that there exists $r>0$ such that
\begin{equation*}\|\ubj (x,-\tau _{\epj })\|_{C^{1+\alpha}\left (
\overline{ \mathcal{B}}_{\epj }\cap \overline{B}_{r} (0)
\right )}\leq C_{r}.\end{equation*}
Moreover, we assume that $\|\bar h_{\epj }(x',t-\tau _{\epj }+\tau )\|_{C^{1}(K)}
+ \|\nabla _{x'}\bar h_{\epj }(x',t-\tau _{\epj }+\tau )\|_{C^{\alpha,\frac{\alpha}{2}}(K)}$
are uniformly bounded for every compact set $K\subset \mathbb{R}^{N-1}\times [-\tau ,+\infty )$.

Then, there holds that
\begin{equation*}\begin{split}
\bar u&\in C^{\alpha , \frac{\alpha }{2}}\bigl (\{x_{1}\geq 0,\,t\geq -\tau \}\bigr ),\,
\nabla \ubar \in C\bigl (\{0\le x_{1}<r,\,t\geq -\tau \}\bigr ),\\
&\ubj (x,-\tau _{\epj })\to \bar u(x,-\tau )\quad \text{uniformly on compact subsets of }\{x_{1}>0\}.
\end{split}\end{equation*}

In any case ($\tau , \g $ be infinite or finite)
\begin{equation*}|\nabla \ubj (0,0)|\to |\nabla \bar u(0,0)|.\end{equation*}
\end{lemma}

\section{Approximation result}
\setcounter{equation}{0}

In this section we  prove that, under certain  assumptions, a
classical supersolution to problem ($P$) is the uniform limit of a
family of supersolutions  to  problem ($\pep$) (Theorem \ref{approx}), and
we prove an analogous result for subsolutions (Theorem \ref{approx2}).
Also, we prove that for compactly supported initial data, limit solutions have bounded
support (Proposition \ref{acotada}).


The following construction follows the lines of Theorem 5.2 in \cite{LVW1}. In our case we have to be more careful with the construction of the initial data.

\begin{theorem} \label{approx}
Let $\wt u$ be a classical supersolution to ($P$) in $Q_T$  with $\wt u\in C^1(\overline{\{\wt u>0\}})$ and such that $ \{\wt u>0\}$ is bounded.
Assume, in addition, that there exist $\delta _{0},s_0>0$ such that
\begin{align*}
& |\nabla \wt{u}^{+}|\le \sqrt {2M(x,t)-\delta _{0}}\quad \text{ on }
Q_T\cap \partial \{\wt{u}>0\},\\
& |\nabla \wt{u}|>\delta_0 \quad \text{ in } Q_T\cap \{0< \wt{u}<s_0 \}.
\end{align*}

Let $w^\ep$ be a solution of the heat equation in $\R^N\times(0,T)$ such that $\frac{\wep(x,t)}\ep\to w_0(x,t)$ uniformly
in $\R^N\times[0,T]$ with $w_0\in C(\R^N\times[0,T])$ and $w_0\ge -1+\delta_1$ for a certain positive constant $\delta_1$.

Then, there
exists a family $\uep \in C(\overline{Q_T})$, with
$\nabla \uep \in L^{2}_{\rm loc}(\overline{Q_T})$,
of weak supersolutions to ($\pep $) in $Q_T$, such that,
as $\ep \to 0$, $\uep \to \wt{u}$ uniformly in $\overline{Q_T}$.

\end{theorem}

\begin{proof}

{\em Step}  I. Construction of the family $\uep $.
Let $0< \theta <\delta_1$ be such that
\begin{equation*}
\int_{1-\theta}^{1}(s+W)f(s)\, ds = \frac{\delta_0}{8},
\end{equation*}
where $W$ is a suitable uniform bound of $\|w^{\ep}/\ep\|_{L^\infty(\{\wt u>0\})}$.
For every $\ep >0$ small, we define the
domain $D^{\ep }= \{ \wt{u}<(1-\theta)\ep\}\subset Q_T$.



Let $\zep $ be the bounded solution to
$$\Delta \zep - \zep_t = (\zep + \wep)f_{\ep}(\zep)\quad \mbox{in } D^{\ep},$$
with boundary data
\begin{align*}
\zep (x,t)&=
\begin{cases}(1-\theta )\ep &\text{on } \partial D^{\ep}\cap {t>0},\\
z_{0}^{\ep}(x) &\text{in }  D^{\ep }\cap \{t=0\}.
\end{cases}
\end{align*}
In order to give the initial data $z_{0}^{\varepsilon }$,
we let $\psi^\ep(s,x)$
be the solution to (\ref{eq4.1}) with
\begin{equation*}a=1-\theta ,\quad
b=\int_{-w^\ep(x,0)/\ep}^{1-\theta}(s+\frac{w^\ep(x,0)}{\ep})f(s)\, ds,\quad \omega_0=\frac{w^\ep(x,0)}\ep.
\end{equation*}
Assume first that $|\nabla \wt u|$ is smooth. Then, by extending $|\nabla \tilde u(x,0)|$
to the whole $\R^N$ as a positive function, we let
\begin{equation*}\varphi^\ep (\xi,x)=\psi^\ep \Bigl (\frac{1-\theta -\xi }{|\nabla {\wt u}(x,0)|},x\Bigr ),
\end{equation*}
and we define
\begin{equation*}z_{0}^{\varepsilon }(x)=\ep {\varphi^\ep }\Bigl (\frac{1}{\ep }{\wt{u}(x,0)},x\Bigr ).
\end{equation*}


If $\wt u$ is not regular enough, we can replace $|\nabla {\wt u}(x,0)|$ by a smooth approximation $F_{\ep}(x)$ so that
the initial datum $z_0^\ep$ is $C^{1+\alpha}$. We leave the details to the reader.


Finally, we define the family $\uep $ as follows:
\begin{equation*}\uep = \begin{cases}\wt{u}&\text{in }\{\wt{u}\geq (1-\theta )\ep \},\\
\zep &\text{in } \overline{D^{\ep }}.
\end{cases}
\end{equation*}

\noindent {\em Step} II. Passage to the limit.
If $(x,0)\in \overline{D^{\ep }}$, we have $0\leq \frac{1}{\ep }{\wt{u}(x,0)}\leq 1-\theta $.
Since, from Lemma 3.1, we know that $-\wep(x,0)/\ep\leq \psi^\ep(s,x)\leq 1-\theta $
for $ s\geq 0$, it follows that $-\wep(x,0)\leq \zep(x,0)\leq (1-\theta )\ep $.
Since $f_{\ep }(s)\ge 0$, constant functions larger than $-{w^\ep(x,t)}$ are supersolutions to
($\pep$). Therefore, $(1-\theta)\ep$ is a supersolution if $\ep<\ep_1$ -- for some $\ep_1>0$ -- and we may apply the comparison principle for bounded super and subsolutions
of ($\pep$) to conclude that $-\wep\leq \zep \leq (1-\theta )\ep $.


Hence,
\begin{equation*}\sup _{\overline{Q_T}}|\uep -\wt{u}|=
\sup_{D^\ep}|\zep-\wt{u}|\leq C\ep \end{equation*} and therefore,
the convergence of the family $\uep $ follows.





\noindent {\em Step} III. Let us show that there exists $\ep _{0}>0$ such
that the functions $\uep $ are supersolutions to ($\pep$) for $\ep <\ep _{0}$.

If $\uep >(1-\theta )\ep $, then $\uep = \wt{u}$, which by hypothesis is supercaloric.
Since $f_{\ep }(s)\ge 0$ and $(1-\theta)\ep\ge-w^\ep$ if $\ep<\ep_1$, it follows that $\uep $ are supersolutions to ($\pep$) here.

If $\uep <(1-\theta )\ep $, then we are in $D^{\ep }$ and therefore, by
construction, $\uep $ are solutions to ($\pep$).

That is, the $\uep $'s are continuous functions, and they are
piecewise  supersolutions to  ($\pep$). In order to see that $\uep $ are
globally  supersolutions to ($\pep$),  it suffices to see that the jumps of
the gradients (which occur at smooth surfaces),  have the right sign.

To this effect, we will show that there exists $\ep _{0}>0$ such that
\begin{equation}
\label{punto}
|\nabla \uep |\geq \sqrt {2M(x,t)-\delta _{0}/2}\quad \text{ on
}\partial\{\wt{u}<(1-\theta )\ep \},\,\text{ for }
\ep <\ep _{0}.
\end{equation}

Assume that (\ref{punto}) does not hold. Then, for every $j\in \mathbb{N}$, there
exist $\ep _{j}>0$ and $(x_{\epj },t_{\epj })\in Q_T$, with
\begin{equation*}\ep _{j}\to 0\quad \text{and} \quad (x_{\epj },t_{\epj })\to (x_{0},t_{0})\in
 \partial \{\wt{u}>0\},\end{equation*}
such that
\begin{equation}
\label{punto2}
\uepj (x_{\epj },t_{\epj })=(1-\theta )\epj \quad \text{ and }\quad
|\nabla \uepj (x_{\epj },t_{\epj })|<\sqrt {2M(x_{\epj},t_{\epj})-\delta _{0}/2}.
\end{equation}

{}From now on we will drop the subscript $j$ when referring to the
sequences defined above and $\ep \to 0$ will mean $j\to \infty $.

We can assume (performing a rotation in the space variables if necessary)
that there exists a family $g_{\ep }$ of
smooth functions such that, in a neighborhood of $(\xe , \te )$,
\begin{equation}\label{5.7}\begin{split}
\{\uep =(1-\theta )\ep \}&=\{(x,t)\,/\,x_{1}-{\xe }_{1}=g_{\ep }(x'-{\xe }',t-\te )\},\\
\{\uep <(1-\theta )\ep \}&=\{(x,t)\,/\,x_{1}-{\xe }_{1}>g_{\ep }(x'-{\xe }',t-\te )\},
\end{split}\end{equation}
where there holds that
\begin{equation*}g_{\ep }(0,0)=0,\quad |{\nabla }_{x'} g_{\ep }(0,0)|\to 0,\quad
\ep \to 0.\end{equation*}
We can assume that (\ref{5.7}) holds in $\bigl ( B_{\rho }(\xe )\times (t_{\ep }-\rho ^{2},
t_{\ep }+\rho ^{2})\bigr )\cap \{0\le t\le T\}$ for
some $\rho >0$.

Let us now define
\begin{equation*}
\ub (x,t)={\frac{1}{\ep }}{\uep (\xe +\ep x,\te +\ep ^{2} t)}, \quad
\bar g_{\ep} (x',t) = {\frac{1}{\ep }}{g_{\ep} (\ep x',\ep ^{2} t)},
\end{equation*}
and let
\begin{equation*}\tau _{\ep }=\frac{\te }{\ep ^{2}}\ ,\  \g _{\ep }=\frac{T-t_{\ep }}{\ep ^{2}}.\end{equation*}
We have, for a subsequence,
\begin{equation*}\tau _{\ep }\to \tau \ ,\ \g _{\ep }\to \g \end{equation*}
where $0\leq \tau ,\g \leq +\infty $ and $\tau $ and $\g $ cannot be both finite.

We now let
\begin{equation*}\mathcal{A}_{\varepsilon}=\left \{ (x,t)\,\,/\,\, |x|<\frac{\rho }{\ep }\, ,
\,\,-\min (\tau _{\ep },\frac{\rho ^{2}}{\ep ^{2}})<t<\min (\g _{\ep },
\frac{\rho ^{2}}{\ep ^{2}})\right \}.
\end{equation*}
Then,  the functions $\ub $ are weak solutions to
\begin{equation*}\begin{array}{rll}
\displaystyle
&\Delta \ub - \ub_t = \left(\ub + \displaystyle\frac{w^{\ep}(\xe + \ep x, \te + \ep^2 t)}{\ep}\right) f(\ub )
& \text{ in }\{x_{1}>\bar g_{\ep }(x',t)\}\cap \mathcal{A}_{\varepsilon},\\
&\displaystyle  \ub =1-\theta  &\text{ on }\{x_{1}=\bar g_{\ep }(x',t)\}\cap
\mathcal{A}_{\varepsilon},\\
  -\hskip-10pt&\displaystyle\frac{w^{\ep}(\xe + \ep x, \te + \ep^2 t)}{\ep}\leq \ub \leq 1-\theta  &
\text{ in }\{x_{1}\geq \bar g_{\ep }(x',t)\}\cap \overline{\mathcal{A}_{\varepsilon}}.
\end{array}\end{equation*}

Note that  we are under the hypotheses
of Lemma \ref{lema4.4}.
Then, there exists a function $\bar u$ such that, for a subsequence,
\begin{gather*}
\bar u\in C^{2+\alpha , 1+ \frac{\alpha }{2}}\bigl (\{x_{1}\geq 0,\,-\tau <t<\g \}\bigr ),\\
\ub \to \bar u\quad \text{uniformly on compact subsets of }
\{x_{1}>0,\,-\tau <t<\g \},\\
\begin{array}{rll}
&\Delta \bar u- \bar u_{t}= (\bar u + w_0(x_0, t_0))f(\bar u)\quad &\text{in
}\{x_{1}>0,\,-\tau <t<\g \},\\
&\bar u=1-\theta \, \quad &\text{on }\{x_{1}=0,\,-\tau <t<\g \},\\
-\hskip-10pt&w_0(x_0, t_0)\leq \bar u\leq 1- \theta \quad &\text{in }\{x_{1}\geq 0,\,-\tau <t<\g \}.
\end{array}
\end{gather*}

We will divide the  remainder of the proof into two cases, depending on
whether $\tau =+\infty $ or $\tau <+\infty $.



\noindent {\em Case} I. Assume $\tau =+\infty $.

In this case, Lemma \ref{lema4.4} also gives
\begin{equation*}|\nabla \ub (0,0)|\to |\nabla \bar u(0,0)|.\end{equation*}

On the other hand, $\bar u$ satisfies the hypotheses of Lemma \ref{lema4.3} and
therefore,
\begin{equation*}|\nabla \bar u| \geq \sqrt {2M(x_0,t_0)-\delta _{0}/4}\quad \text{on }
\{x_{1}=0\},
\end{equation*}
which yields
\begin{equation*}|\nabla \ub (0,0)|\geq \sqrt {2M(x_0,t_0)- 3{\delta _{0}}/{8}},
\end{equation*}
for $\ep $  small. But this gives
\begin{equation*}|\nabla \uep (\xe ,\te )|\geq \sqrt {2M(\xe,\te)-{\delta _{0}}/{2}},
\end{equation*}
for $\ep $ small. This contradicts (\ref{punto2}) and completes the proof in case
$\tau =+\infty $.

\noindent {\em Case} II. Assume $\tau <+\infty $. (In this case $\g =+\infty$.)

There holds that $\ub (x,-\tau _{\ep })=\frac{1}{\ep }\uep (\xe +\ep x, 0)$, then
\begin{equation}\label{datoinicial}
\ub (x,-\tau _{\ep })=\varphi^\ep \Bigl (\frac{1}{\ep}\wt{u}(\xe +\ep x,0),\xe +\ep x \Bigr ).
\end{equation}

Here we want to apply the result of Lemma \ref{lema4.4} corresponding to $\tau <+\infty $.
In fact, we can see that there exist $C, r >0$ such that
$\|\ub (\cdot,-\tau _{\ep })\|_{C^{1+\alpha}\left (\overline{B}_{r}(0)\right )}\le C$.

Now Lemma \ref{lema4.4} gives, for a subsequence,
\begin{equation*}\begin{split}
\bar u&\in C^{\alpha , \frac{\alpha }{2}}\bigl (\{x_{1}\geq 0,\,t\geq -\tau \}\bigr ),\\
\ub (x,-\tau _{\ep })&\to \bar u(x,-\tau )\quad \text{uniformly on compact subsets of }\{x_{1}>0\}.
\end{split}\end{equation*}
Therefore, we get that (recall that in the case we are considering $t_0=0$),
\begin{equation*}\bar u(x,-\tau )=\bar \varphi \Bigl (1- \theta -|\nabla \wt{u}^{+}(x_{0},t_{0})|\,x_{1}, x_0\Bigr ).
\end{equation*}
where $\bar\varphi(s,x)=\psi\Big(\displaystyle\frac{1-\theta-s}{|\nabla\wt u(x,0)|},x\Big)$ and $\psi(s,x)$ is the solution of  (\ref{eq4.1}) with
\begin{equation*}a=1-\theta ,\quad
b=\int_{-w_0(x,0)}^{1-\theta}(s+{w_0(x,0)})f(s)\, ds,\quad \omega_0={w_0(x,0)}.
\end{equation*}


Thus,
\begin{equation*}\bar u(x,-\tau )=\psi (x_{1},x_0).\end{equation*}

Since the function
$\psi (x_{1},x_0)$
is a stationary  solution to equation ($P_{0}$), bounded for $x_{1}\geq 0$,
and $\bar u= \psi $ on the parabolic boundary of the domain
$\bigl \{x_{1}> 0,\, t>-\tau \bigr \}$,
we conclude that
\begin{equation*}\bar u(x,t)=\psi (x_{1},x_0)
\quad \text{ in }
\bigl \{x_{1}\geq 0,\, t\geq -\tau \bigr \}.
\end{equation*}
It follows from Lemma 3.1 and the choice of $\theta$ that
\begin{equation*}
\displaystyle\frac12|\nabla \bar u(0,0)|^2=\frac12\big(\psi_s(0,x_0)\big)^2=\int_{-w_0(x_0,t_0)}^{1-\theta}
(s+w_0(x_0,t_0))f(s)\,ds\ge M(x_0,t_0)-\frac{\delta_0}8.
\end{equation*}
This is,
\begin{equation*}|\nabla \bar u|\ge \sqrt {2M(x_0,t_0)-\delta _{0}/4}\quad \text{on } \{ x_{1}=0, t\geq -\tau \}.
\end{equation*}
But Lemma \ref{lema4.4} gives
\begin{equation*}|\nabla \ub (0,0)|\to |\nabla \bar u(0,0)|,
\end{equation*}
which yields
\begin{equation*}|\nabla \ub (0,0)|\geq \sqrt {2M(x_0,t_0)- 3 {\delta _{0}}/{8}},
\end{equation*}
for $\ep $  small. Then,
\begin{equation*}|\nabla \uep (\xe ,\te )|\geq \sqrt {2M(\xe,\te)-{\delta _{0}}/{2}},
\end{equation*}
for $\ep $ small. This contradicts (\ref{punto2}) and completes the proof in case
$\tau <+\infty $.  \end{proof}

\begin{rem} Observe that from the construction of $\uep$ done in the previous proof, it follows that
$$\uep \equiv \wt{u} \qquad \text{in } \{ \wt{u}>(1-\theta)\ep \}.$$
\end{rem}

\medskip



\begin{theorem} \label{approx2}
Let $\wt u$ be a classical subsolution to ($P$) in $Q_T$ with $\wt u\in C^1(\overline{\{\wt u>0\}})$  such that $ \{\wt u>0\}$ is bounded.
Assume, in addition, that there exist $\delta _{0}>0$ such that
\begin{equation*} |\nabla \wt{u}^{+}|\ge \sqrt {2M(x,t)+\delta _{0}}\quad \text{ on }
Q_T\cap \partial \{\wt{u}>0\}.\end{equation*}

Let $w^\ep$ be a solution of the heat equation in $\R^N\times(0,T)$ such that $\frac{\wep(x,t)}\ep\to w_0(x,t)$ uniformly
in $\R^N\times[0,T]$. And assume, moreover that $w_0\in C(\R^N\times[0,T])$ and $w_0(x,t)\ge -1+\delta_1$ for a certain positive constant $\delta_1$.


Then, there
exists a family $\uep \in C(\overline{Q_T})$, with
$\nabla \uep \in L^{2}_{\rm loc}(\overline{Q_T})$,
of weak subsolutions to ($\pep $) in $Q_T$, such that,
as $\ep \to 0$, $\uep \to \wt{u}$ uniformly in $\overline{Q_T}$.
\end{theorem}

\begin{proof}
The proof is analogous to Theorem \ref{approx}. See \cite{LVW1} for a similar result in the case
$\wep=0$.
\end{proof}

Finally, we end this Section by showing that, for compactly supported initial data, the support
of a limit solution of problem ($P$) is bounded.
\begin{proposition}\label{acotada}
Let $u_0\in C(\R^N)$ with compact support. Let $u_0^\ep$ converge uniformly
to $u_0$ with supports converging to the support of $u_0$ and let $w^\ep$ be a solution of the heat equation in $\R^N\times(0,T)$ such that $\frac{\wep(x,t)}\ep\to w_0(x,t)$ uniformly
in $\R^N\times[0,T]$. And assume, moreover that $w_0\in C(\R^N\times[0,T])$ and $w_0(x,t)\ge -1+\delta_1$ for a certain positive constant $\delta_1$. Finally, let $\uep$ be the solution to ($\pep$)
with function $\wep$ and initial condition $u_0^\ep$.

Let $u=\lim\uepj$. Then  $\{ u>0\}$ is bounded. Moreover, $u$ vanishes in finite
time.
\end{proposition}

\begin{proof}
% NOTE (2026 review): as written, omega_0 depends on (x,t) and eps. What is needed is
% a constant -1 < omega_0 < -1+delta_1 (so omega_0 < w_0 everywhere, and the gap
% M - M_{omega_0} is bounded below uniformly).
Let $-1<\omega_0 < \wep(x,t)/\ep$. Then it is easy to check that
\begin{equation}
\label{MM}
M_{\omega_0} = \int_{-\omega_0}^1 (s+\omega_0)f(s)\, ds < M(x,t) =
\int_{-w_0(x,t)}^1 (s+w_0(x,t))f(s)\, ds.
\end{equation}

Let us now consider the following self-similar function
$$V(x,t;T) = (T-t)^{1/2} h(|x|(T-t)^{-1/2}),$$
where $h=h(r)$ is a solution of
\begin{equation}
\label{perfil}
\begin{split}
& h'' + \left(\frac{N-1}{r} + \frac12 r\right)h' + \frac12 h = 0, \quad 0<r<R,\\
& h'(0) = 0,\quad h(r)>0, \quad 0\le r < R,\\
& h(R) = 0, \quad h'(R) = -\sqrt{2M_{\omega_0}}.
\end{split}
\end{equation}
It is proved in \cite{CV}, Proposition 1.1, that there exists a unique $R>0$ and a unique $h$ solution
of (\ref{perfil}).

 Moreover, it can be checked that if one picks $T$ sufficiently large, then
$$V(x,0;T) \ge u_0 + 1 \quad \mbox{in } \{ u_0 > 0 \},$$
and so $V(x,t;T)$ is a strict classical supersolution of ($P$) with bounded support and positive
gradient near its free boundary.

Now, let $\uepj$ be solutions to ($P_{\epj}$) -- with initial data $u_{0}^{\epj} $ converging uniformly to $u_0$ such that
support $u_{0}^{\epj} \to $ support $u_0$ -- such that $u = \lim \uepj$.

By Theorem \ref{approx}, there exists a family $\vepj$ of supersolutions of ($P_{\epj}$) such that
$\vepj\to V$ uniformly on compact sets, and $\vepj(x,0)\ge \uepj(x,0)$.
Therefore, by the comparison principle, we obtain $\uepj\le \vepj$ and passing to the limit
$u(x,t)\le V(x,t;T)$, and the result follows.
\end{proof}



\section{Uniqueness of the limit solution}
\setcounter{equation}{0}

In this section  we arrive at the main point of the article: we prove that, under certain assumptions,
there exists at most one limit solution  to the initial and boundary value problem associated to ($P$) as
long as condition (\ref{condicion_inicial}) is satisfied.


Let us begin with the following Proposition that is the key ingredient in the proof of our main result.

\begin{proposition}
\label{keyprop}
Let $\wt{u}$ be a strict classical  supersolution to ($P$) with bounded support in $\R^N\times(0,T)$ such that
there exists $s_0>0$ so that
$|\nabla \wt{u}|>0$ in $\{0<\wt u<s_0\}$ and let $\wep/\ep$ be solutions to the heat equation in $\R^N\times(0,T)$
converging to $w_0$ uniformly with $w_0\in C(\R^N\times[0,T])$ and $w_0\ge-1+\delta_1$ for a certain
positive constant $\delta_1$.

Let $\uep$ be
solutions to ($\pep$) with function $\wep$ and initial condition $u_0^\ep$, where $u_{0}^\ep$ are uniform approximations of $u_0$ with
$\mbox{\rm support } u_{0}^\ep \to \mbox{\rm support } u_0$. Then
\begin{equation*}
\limsup_{\ep\to 0+} \uep(x,t)\le \wt{u}(x,t)
\end{equation*}
for every $(x,t)\in Q_T$.
\end{proposition}
\begin{proof}
Let $\wt{u}$ be a strict classical  supersolution of ($P$). Let us first, define the following regularization
$$u(x,t) = (\wt{u}(x,t+h)-\eta)^+,$$
for $h,\eta>0$ small. So  that $u$ is a  strict classical  supersolution of ($P$) with $C^1_x$
free boundary, $C^1(\overline{\{u>0\}})$  and $|\nabla u|>\delta_0>0$ in a neighborhood of its free boundary.
So, by Theorem \ref{approx}, there exists $\vep$ supersolution of ($\pep$) such that
$\vep\to u$ uniformly in $Q_{T-h}$.

Now, using the comparison principle, we conclude that $\uep\le \vep$ in $Q_{T-h}$, and the
Proposition now follows letting first $\ep\to 0+$ and then $h, \eta\to 0+$.
\end{proof}


Finally, we arrive at the main point of the paper: The uniqueness of limit solutions of ($P$).

\begin{theorem}
\label{unicidad}
Let the initial datum $u_0$ be Lipschitz, with compact support and satisfy the condition \eqref{S}.
Then there exists at most one limit solution such that its gradient does not vanish near its free boundary as long as
the function $\wep$ in problem $(\pep)$ satisfies condition (\ref{def_w_0}).



More precisely, let $u_0^{\ep_j},\tilde u_{0}^{\ep_k}$ be uniformly Lipschitz continuous in $\R^N$ with uniformly bounded Lipschitz norms and $\ep_j,\,\ep_k\to0$.
Assume that $u_0^{\ep_j}\in C^1(\overline{\{u^{\ep_j}_0>0\}})$, $\tilde u_{0}^{\ep_k}\in C^1(\overline{\{\tilde u^{\ep_k}_0>0\}})$,
$u_0^{\ep_j},\tilde u_{0}^{\ep_k}\to u_0$ uniformly and
${\rm support}\ u_0^{\ep_j}$, ${\rm support}\ \tilde u_{0}^{\ep_k}$$\to {\rm support}\ u_0$. Let
${w^{\ep_j}}/{\ep_j}$ and ${\tilde w^{\ep_k}}/{\ep_k}$
be solutions of the heat equation converging to the same function $w_0\in C(\overline{Q_T})$, uniformly bounded
from below by $-1+\delta_1$ for a certain positive constant $\delta_1$. Also, assume that  $w_0$ satisfies the monotonicity condition
(\ref{mono_w}).


Let $u^{\ep_j}$ (resp. ${\tilde u}^{\ep_k}$) be the solution to $(P_{\ep_j})$ with function $w^{\ep_j}$ and initial datum
$u_{0}^{\ep_j}$ (resp. solution to $(P_{\ep_k})$ with function $\tilde w^{\ep_k}$ and initial datum $ \tilde u_{0}^{\ep_k}$). Let
$u=\lim \uepj$ and $\tilde u=\lim \tilde u^{\ep_k}$. If there exists $s_0>0$ such that $|\nabla \tilde u|>0$ in $\{0<\tilde u<s_0\}$.

Then, $u\le\tilde u$.


\end{theorem}

\begin{proof}
Since  $\tilde u$ is a classical supersolution of ($P$), $\tilde u\in C^1(\{\tilde u>0\})$ and, by Proposition \ref{acotada}, its  support is bounded,  the function
 $\tilde u_\lambda$  as defined in (\ref{rescalada}) satisfies the hypotheses of Proposition \ref{keyprop} in $Q_{T/\lambda^2}\supset Q_T$.  So by
 letting $\lambda\to 1-$ we arrive at
\begin{equation}
\label{desig1}
u(x,t)\le \tilde u(x,t).
\end{equation}

This finishes the proof.



\end{proof}



\begin{theorem}
\label{unicidad2}
Let the initial datum $u_0$ be as in Theorem \ref{unicidad}. Assume that there exists a classical
solution $v$ to ($P$) with initial datum $u_0$ and let
$u_0^{\ep_j}$ be uniformly Lipschitz continuous in $\R^N$ with $\ep_j\to0$,
such that $u_0^{\ep_j}\in C^1(\overline{\{u^{\ep_j}_0>0\}})$,
$u_0^{\ep_j}\to u_0$ uniformly and
${\rm support}\ u_0^{\ep_j}$$\to {\rm support}\ u_0$. Assume
${w^{\ep_j}}/{\ep_j}$ is a  solution of the heat equation converging to $w_0$ uniformly with
$w_0\in C(\R^N\times[0,T])$ and $w_0\ge -1+\delta_1$
in $\R^N\times(0,T)$ for a certain $\delta_1>0$. Also, assume that  $w_0$ satisfies the monotonicity condition
(\ref{mono_w}).

Let  $u^{\ep_j}$ be the solution to $(P_{\ep_j})$ with function $w^{\ep_j}$ and initial datum
$u_{0}^{\ep_j}$ and let $u=\lim \uepj$.
Then, $u=v$.

 In particular, there exists at most one classical solution to ($P$).
\end{theorem}


\begin{proof}Since $u$ is a classical supersolution to ($P$) and $v$ is a classical subsolution, Lemma
\ref{compar} applies and we get that $v\le u$.

On the other hand, if we define $v_{\lambda}$ as in (\ref{rescalada}), with $0<\lambda<\lambda'<1$,
we have that $v_{\lambda}$ is a strict classical supersolution.
% NOTE (2026 review): Proposition \ref{acotada} is stated for limit solutions, not
% for an arbitrary classical solution v. Proving that v has bounded support by comparing
% it with the self-similar supersolution V via Lemma \ref{compar} would be circular,
% since that lemma assumes the subsolution already has bounded support. Either add
% bounded support of v as a hypothesis or justify it separately.
Since $v_\lambda$ has compact
support (see Proposition \ref{acotada}) it
satisfies the hypotheses of Proposition \ref{keyprop}. Thus,
$$u=\lim \uepj \le  v_\lambda.$$

Letting $\lambda\to1^-$ we obtain the desired result.
\end{proof}





\section{Conclusions}

In this paper we have proved that the limits of sequences of solutions to ($\pep$) with different
constitutive functions $\wep$ and initial data $\uep_0$ coincide -- as long as certain monotonicity
assumptions are made -- if the limit of $\wep/\ep$ and of $\uep_0$ are prescribed.

The monotonicity assumptions are necessary to provide strict classical supersolutions as close as
we want to any classical supersolution. This kind of condition was also used with the same purpose
-- in the case in which $\wep=0$ -- in \cite{CP} and \cite{LVW1}. In the latter, a different geometry was considered namely, the domain was a
cylinder, Neumann boundary conditions were given on the boundary of the cylinder and monotonicity in the direction of the cylinder axis was assumed. In \cite{LVW1} it was proved that, if a
classical solution exists and $\wep=0$, then it is equal to any limit of solutions to ($\pep$).

% NOTE (2026 review): the cylinder case below is claimed but not proved in this paper.
In our case, this is with $\wep\neq0$  satisfying (\ref{def_w_0}) and nondecreasing in the direction
of the cylinder axis, the uniqueness result in the presence of a classical solution still holds.

The cylindrical geometry has the advantage of giving the condition of nonvanishing gradient in the
positivity set of any limit solution.  Since in dimension 2 one can prove that limit solutions are classical supersolutions up to the fixed boundary, the uniqueness of limit solutions follows in this case without further assumptions.

\bibliographystyle{amsalpha}
\begin{thebibliography}{xx}


\bibitem{BuL} J.D. Buckmaster and G.S.S. Ludford, {\em Theory of Laminar Flames},
Cambridge University Press, Cambridge, 1982.


\bibitem{CLW1} L.A. Caffarelli, C. Lederman, and N. Wolanski, {\em Uniform estimates and limits
for a two phase parabolic singular perturbation problem}, Indiana Univ. Math. J. {\bf 46 (2)} (1997),
453--490.

\bibitem{CLW2} \bysame , {\em Pointwise and viscosity solutions for the limit of a two phase parabolic
singular perturbation problem}, Indiana Univ. Math. J. {\bf 46 (3)} (1997), 719--740.


\bibitem{CV} L.A. Caffarelli and J.L. Vazquez, {\em A free boundary problem for the heat equation
arising in flame propagation}, Trans. Amer. Math. Soc. {\bf 347} (1995), 411--441. \MR{95e:35097}

\bibitem{FBW} J. Fern\'andez Bonder and N. Wolanski, {\em A free boundary problem in combustion theory}.
Interfaces Free Bound., {\bf 2} (2000), 381--411.

\bibitem{L} A. Langlois, {\em Sur l'\`etude asymptotique d'un syst\`eme parabolique mod\`elisant
des flames presque \`equidiffusives}. PhD dissertation, Ecole Central de Lyon, December 2000.

\bibitem{LVW1} C. Lederman, J.L. Vazquez,  and N. Wolanski, {\em Uniqueness of solution to a free
boundary problem from combustion}. Trans. Amer. Math. Soc. {\bf 353} (2001), no. 2, 655--692.

\bibitem{LVW2} \bysame , {\em Uniqueness of solution to a two-phase free boundary problem from combustion},
 Advances Differential Equations {\bf 6} (2001), no. 12, 1409--1442.

\bibitem{CP} A. Petrosyan, {\em On existence and uniqueness in a free boundary
problem from combustion}, Comm. Partial Differential Equations {\bf 27} (2002), no. 3-4, 763--789.

\end{thebibliography}

\end{document}
